\begin{abstract}
Domain adaptation across distant modalities exposes which inductive biases genuinely transfer and which remain tied to source-domain artifacts. Over five prior workshops, Group~8 built machine learning pipelines for autoregressive next-character prediction on the \texttt{flwrlabs/shakespeare} corpus (4.22~million rows, 76 character classes, 791,902:1 imbalance), culminating in classical \texttt{scikit-learn} baselines (Workshop~1), a 4-layer Pre-LN \texttt{CharTransformer} with post-hoc logit adjustment (Workshop~2), a Double+Dueling DQN text-control agent (Workshop~3), a 1D residual convolutional VAE for anomaly detection (Workshop~4), and a 1D WGAN-GP with Gumbel-Softmax straight-through sampling (Workshop~5). In this final workshop, we audit every component from Workshops~1--5 and adapt viable techniques to dermatoscopic skin lesion diagnosis on the public \textbf{HAM10000 / ISIC~2018 Task~3} dataset (10,015 images across seven lesion classes, 58.3:1 class imbalance, linked clinical metadata including patient age, biological sex, anatomical localization, and diagnostic verification method, plus pixel-level lesion segmentation masks).

We deliberately separate positive architectural and methodological transfers from honest negative results. Leakage-free entity-grouped splitting (\texttt{GroupShuffleSplit} by \texttt{lesion\_id}), Pre-LN Transformer blocks, post-hoc logit adjustment, 2D residual VAE anomaly scoring with Free Bits and cyclical KL annealing, and FP32 WGAN-GP gradient penalties adapt cleanly to 2D medical imaging. Conversely, character embeddings, causal attention masks, Gumbel-Softmax discretization, lexical n-gram features, and the Workshop~3 DQN agent do not transfer to static single-step diagnostic imaging. Furthermore, freezing Workshop~2 Shakespeare Transformer blocks inside a hybrid CNN-Transformer (\texttt{MedHybridTransformer}) serves as an explicit empirical control testing extreme cross-domain feature reuse against training from scratch and ImageNet-pretrained \texttt{ResNet-18}.

On the untouched lesion-grouped test set (\WSixTestImgs{} images from \WSixTestLesions{} lesions), our best adapted classifier (\WSixBestModel{}) achieves a 7-class balanced accuracy of \WSixBestBalAcc{}, malignancy AUC-ROC of \WSixBestMalAUC{}, PR-AUC of \WSixBestMalPRAUC{}, and post-calibration expected calibration error of \WSixBestECE{}. At a validation-selected high-sensitivity screening operating point, the system reaches \WSixBestSensScr{} sensitivity and \WSixBestSpecScr{} specificity, yielding positive clinical net benefit on Decision Curve Analysis over biopsy-all and biopsy-none strategies. Quantitative interpretability verification against Tschandl lesion masks measures \WSixXaiLesionEnergy{} saliency energy inside lesions and \WSixXaiPointingAcc{} pointing-game accuracy across 15~TP/FN/FP test cases, while exposing five concrete clinical failure modes including border-shortcut attention and subgroup disparities across patient age and anatomical site. We conclude that our models demonstrate credible decision-support potential for primary-care teledermatology triage only under human-in-the-loop oversight, external multi-center validation across diverse skin phototypes, formal FDA De~Novo or 510(k) clearance with a Predetermined Change Control Plan, and continuous post-market drift monitoring.
\end{abstract}
